% from ega3-1-fr.tex
\oldpage[0\textsubscript{III}]{27}
En effet, d'après les conditions a), b), c) de
(\hyperref[0.11.1.1-fr]{11.1.1}), dire que
$Z_{r+1}(E_2^{pq})=Z_r(E_2^{pq})$ équivaut à dire que $d_r^{pq}=0$, et dire
que
$B_r(E_2^{p+r,q-r+1})=B_{r+1}(E_2^{p+r,q-r+1})$ équivaut aussi à dire que
$d_r^{pq}=0$~; le lemme résulte aussitôt de cette remarque.
\end{proof}
\subsection{La suite spectrale d'un complexe filtré}
\label{subsection:0.11.2-fr}

\begin{env}[11.2.1]
\label{0.11.2.1-fr}
Étant donnée une catégorie abélienne $\mathcal C$, nous conviendrons de
désigner par des notations telles que $K^\bullet$ les \emph{complexes}
$(K^i)_{i\in\mathbb Z}$ d'objets de $\mathcal C$ dans lesquels l'opérateur de
dérivation est de degré $+1$, par des notations telles que $K_\bullet$ les
complexes $(K_i)_{i\in\mathbb Z}$ d'objets de $\mathcal C$ dans lesquels
l'opérateur de dérivation est de degré $-1$. À tout complexe
$K^\bullet=(K^i)$ dont l'opérateur de dérivation $d$ est de degré $+1$, on
peut associer un complexe $K'_\bullet=(K'_i)$ en posant $K'_i=K^{-i}$,
l'opérateur de dérivation $K'_i\to K'_{i-1}$ étant l'opérateur
$d:K^{-i}\to K^{-i+1}$~; et \emph{vice versa}, ce qui permettra suivant les
circonstances de considérer l'un ou l'autre type de complexes et de traduire
tout résultat pour l'un des types en résultats pour l'autre. Nous désignerons
de même par des notations telles que $K^{\bullet\bullet}=(K^{ij})$ (resp.
$K_{\bullet\bullet}=(K_{ij})$) des \emph{bicomplexes} d'objets de
$\mathcal C$ dans lesquels les \emph{deux} opérateurs de dérivation sont de
degré $+1$ (resp. $-1$)~; on passe encore de l'un à l'autre type en changeant
les signes des indices, et on a des notations et remarques analogues pour des
multicomplexes quelconques. La notation $K^\bullet$ ou $K_\bullet$ sera aussi
utilisée pour des \emph{objets gradués} de $\mathcal C$, de type $\mathbb Z$,
qui ne sont pas nécessairement des complexes (ou que l'on peut considérer
comme tels pour des opérateurs de dérivation nuls)~; par exemple, nous
écrirons
$\mathrm H^\bullet(K^\bullet)=(\mathrm H^i(K^\bullet))_{i\in\mathbb Z}$ la
\emph{cohomologie} d'un complexe $K^\bullet$ dont l'opérateur de dérivation
est de degré $+1$, par
$\mathrm H_\bullet(K_\bullet)=(\mathrm H_i(K_\bullet))_{i\in\mathbb Z}$
l'\emph{homologie} d'un complexe $K_\bullet$ dont l'opérateur de dérivation
est de degré $-1$~; quand on passe de $K^\bullet$ à $K'_\bullet$ par
l'opération décrite ci-dessus, on a
$\mathrm H_i(K'_\bullet)=\mathrm H^{-i}(K^\bullet)$.

Rappelons à ce propos que pour un complexe $K^\bullet$ (resp. $K_\bullet$),
nous écrirons en général
$Z^i(K^\bullet)=\operatorname{Ker}(K^i\to K^{i+1})$ («~objet des cocycles~»)
et $B^i(K^\bullet)=\operatorname{Im}(K^{i-1}\to K^i)$ («~objet des
cobords~») (resp.
$Z_i(K_\bullet)=\operatorname{Ker}(K_i\to K_{i-1})$ («~objet des cycles~»)
et $B_i(K_\bullet)=\operatorname{Im}(K_{i+1}\to K_i)$ («~objet des
bords~»)) de sorte que
$\mathrm H^i(K^\bullet)=Z^i(K^\bullet)/B^i(K^\bullet)$ (resp.
$\mathrm H_i(K_\bullet)=Z_i(K_\bullet)/B_i(K_\bullet)$).

Si $K^\bullet=(K^i)$ (resp. $K_\bullet=(K_i)$) est un complexe dans
$\mathcal C$, et $T:\mathcal C\to\mathcal C'$ un foncteur de $\mathcal C$
dans une catégorie abélienne $\mathcal C'$, nous désignerons par
$T(K^\bullet)$ (resp. $T(K_\bullet)$) le complexe $(T(K^i))$ (resp.
$(T(K_i))$) dans $\mathcal C'$.

Nous ne revenons pas sur la définition des $\partial$-foncteurs (T, 2.1),
sauf pour signaler que nous dirons aussi $\partial$-foncteur au lieu de
$\partial^*$-foncteur lorsque le morphisme $\partial$ diminue le degré d'une
unité, le contexte devant chaque fois préciser ce point s'il peut y avoir
doute.

Enfin, nous dirons qu'un \emph{objet gradué} $(A_i)_{i\in\mathbb Z}$ de
$\mathcal C$ est \emph{limité inférieurement} (resp. \emph{supérieurement})
s'il existe $i_0$ tel que $A_i=0$ pour $i<i_0$ (resp. $i>i_0$).
\end{env}

\begin{env}[11.2.2]
\label{0.11.2.2-fr}
Soit $K^\bullet$ un complexe de $\mathcal C$ dont l'opérateur de dérivation
$d$ est de degré $+1$, et supposons-le muni d'une \emph{filtration}
$F(K^\bullet)=(F^p(K^\bullet))_{p\in\mathbb Z}$ formé de sous-
